\wprowadzenie

% Wprowadzenie ma cztery zadania, w tej kolejności: osadzić problem, wskazać
% lukę, postawić cel, zapowiedzieć strukturę. Objętość: 3-5 stron.
% Usuń akapity instruktażowe w miarę pisania własnego tekstu.

\wyroznienie{Tło problemu.} Zacznij od problemu, który istnieje niezależnie od Twojego pakietu, i pokaż, dlaczego jest ważny dla badaczy nauk społecznych. Nie zaczynaj od narzędzia – zaczynaj od pytania, na które badacz nie potrafi dziś odpowiedzieć. Każde twierdzenie o stanie rzeczy opatrz cytowaniem w postaci \kod{\textbackslash parencite\{klucz\}}.

\wyroznienie{Luka.} Wyjaśnij, czego brakuje: metoda jest opisana w literaturze, ale nie ma jej implementacji; istniejące implementacje wymagają przygotowania programistycznego; dostępne narzędzia nie obsługują danych określonego rodzaju. Luka musi być sprawdzalna – wymień narzędzia, które przejrzałeś, i napisz, czego w nich nie ma.

\wyroznienie{Cel i pytanie.} Sformułuj cel jednym zdaniem, w formie czynności: zaprojektowanie, implementacja i weryfikacja pakietu języka R udostępniającego procedurę X. Następnie postaw pytanie, na które praca odpowiada. Pytania zaczynające się od „jak" i „dlaczego" prowadzą do analizy; pytania „co" i „ile" prowadzą do opisu.

\wyroznienie{Materiał i metoda.} Krótko: na jakich danych sprawdzasz narzędzie (dane wygenerowane własną procedurą, zbiór otwarty, jedno i drugie) i jak weryfikujesz poprawność (przypadki o znanym wyniku, testy własności, porównanie z wynikami z artykułu źródłowego).

\wyroznienie{Zakres i ograniczenia.} Napisz, czego praca nie obejmuje. Ograniczenie postawione samodzielnie jest mocniejsze niż zarzut recenzenta.

\wyroznienie{Struktura.} Jedno zdanie o każdym rozdziale.

% Wzorce zapisu, przydatne w całej pracy:
%
%   cytowanie zwykłe            \parencite{klucz}
%   cytowanie ze stroną         \parencite[s.~15-21]{klucz}
%   cytowanie w zdaniu          \textcite{klucz} pokazuje, że...
%   kilka pozycji naraz         \parencite{klucz1,klucz2}
%   nazwisko do indeksu         \osoba{Nazwisko}{Imię}
%   nazwa funkcji               \fun{przygotuj_dane}
%   nazwa argumentu             \argument{skala}
%   nazwa pakietu               \pkg{ggplot2}
%   ścieżka lub plik            \plik{R/algorytm.R}
%   wprowadzany termin          \termin{fuzzyfikacja}
%   wyróżnienie treściowe       \wyroznienie{tekst}
